\documentclass{article}

\usepackage{amsmath,amssymb}
\usepackage{xurl}
\usepackage[hidelinks]{hyperref}
\setlength{\emergencystretch}{2em}

% Shared commands for notes in docs/paper-gaps/.

\DeclareUnicodeCharacter{2080}{\ifmmode{}_{0}\else\textsubscript{0}\fi}
\DeclareUnicodeCharacter{2081}{\ifmmode{}_{1}\else\textsubscript{1}\fi}
\DeclareUnicodeCharacter{2082}{\ifmmode{}_{2}\else\textsubscript{2}\fi}
\DeclareUnicodeCharacter{2099}{\ifmmode{}_{n}\else\textsubscript{n}\fi}

\newcommand{\Lean}{\textsc{Lean}}
\ExplSyntaxOn
\NewDocumentCommand{\leanid}{m}{
  \ifmmode
    \textsf{\detokenize{#1}}
  \else
    \group_begin:
    \urlstyle{sf}
    \tl_set:Nn \l_tmpa_tl {#1}
    \tl_replace_all:Nnn \l_tmpa_tl {\_} {_}
    \exp_args:NV \path \l_tmpa_tl
    \group_end:
  \fi
}
\ExplSyntaxOff
\providecommand{\tr}{\operatorname{tr}}
\newcommand{\tnleanIssueBase}{https://github.com/LionSR/TNLean/issues/}
\newcommand{\tnleanPrBase}{https://github.com/LionSR/TNLean/pull/}
\newcommand{\tnleanDocBase}{https://sirui-lu.com/TNLean/docs/}
\newcommand{\ghissue}[1]{\href{\tnleanIssueBase#1}{GitHub issue \##1}}
\newcommand{\ghpr}[1]{\href{\tnleanPrBase#1}{PR \##1}}
\newcommand{\doclink}[1]{\href{\tnleanDocBase#1.html}{\path{#1}}}

% Dirac notation and the identity operator, as in blueprint/src/macros/common.tex.
% \providecommand keeps note-local definitions authoritative.
\usepackage{dsfont}
\providecommand{\ket}[1]{|#1\rangle}
\providecommand{\bra}[1]{\langle #1|}
\providecommand{\braket}[2]{\langle #1 | #2 \rangle}
\providecommand{\ketbra}[2]{|#1\rangle\!\langle #2|}
\providecommand{\Id}{\mathds{1}}
\providecommand{\one}{\mathds{1}}
\providecommand{\C}{\mathbb{C}}
\providecommand{\MN}[1]{M_{#1}(\C)}
\providecommand{\MD}{M_D(\C)}

% External referencing for paper sources.  The build helper writes sanitized
% auxiliary files under build/paper-aux/ and excludes duplicated source labels.
\usepackage{xr}

% Import a generated paper auxiliary file with a caller-supplied label prefix.
% The candidate paths support builds from docs/paper-gaps/, the repository root,
% and build/paper-aux/ itself.
\newcommand{\importpaperaux}[2]{%
  \IfFileExists{../../build/paper-aux/#2.aux}{%
    \externaldocument[#1]{../../build/paper-aux/#2}%
  }{%
    \IfFileExists{build/paper-aux/#2.aux}{%
      \externaldocument[#1]{build/paper-aux/#2}%
    }{%
      \IfFileExists{#2.aux}{%
        \externaldocument[#1]{#2}%
      }{}%
    }%
  }%
}

% General external reference macros
\newcommand{\paperref}[2]{\ref{#1-#2}}
\newcommand{\papereqref}[2]{(\ref{#1-#2})}

% CPSV16 source (1606.00608 / MPDO-22-12-17-2.tex) external document and shortcuts
\importpaperaux{cpsv16-}{1606.00608/MPDO-22-12-17-2-tnlean}
\newcommand{\cpsvref}[1]{\paperref{cpsv16}{#1}}
\newcommand{\cpsveqref}[1]{\papereqref{cpsv16}{#1}}

% Verdict marker: \gapnote{<kind>}{<status>}. Kind is the policy
% classification (clarification, local-correction, scope-restriction,
% unfaithful, false-source, open-gap); status is open, wip (elimination
% underway), resolved, or historical. Machine-read by the site index;
% prints a verdict line.
\newcommand{\gapnote}[2]{\par\noindent\textbf{Verdict:} #1 (#2).\par}


\title{The Projector-Closure Sufficient Condition for Canonical Form}
\author{}
\date{2026-07-10}

\begin{document}
\maketitle
\gapnote{clarification}{resolved}

\section{The source assertion}

Cirac, P\'erez-Garc\'ia, Schuch, and Verstraete give a sufficient condition for
canonical form in Section~II of their analysis of matrix product density
operators \cite{Cirac2016MPDO_arXiv}. Assume that the tensor has no nontrivial
periodic vectors. The projector-closure hypothesis states that every
orthogonal projection $P$ satisfying
\begin{align}
    PA^i
        &=PA^iP,
        \qquad\text{for every }i,
    \label{eq:cpsv16_projector_closure_canonical_form_source_invariant_projection}
\end{align}
also satisfies
\begin{align}
    A^iP
        &=PA^iP,
        \qquad\text{for every }i.
    \label{eq:cpsv16_projector_closure_canonical_form_source_coinvariant_projection}
\end{align}
The paper asserts that these conditions put the tensor in canonical form
\cite[lines~253--254]{Cirac2016MPDO_arXiv}.

\section{The algebraic projection step}

The projector hypothesis has a normalization-free consequence. Suppose that
$P$ is an orthogonal projection such that
\begin{align}
    (I-P)A^iP
        &=0,
        \qquad\text{for every }i.
    \label{eq:cpsv16_projector_closure_canonical_form_lower_corner_zero}
\end{align}
Then
\begin{align}
    (I-P)A^i
        &=(I-P)A^iP+(I-P)A^i(I-P)
    \label{eq:cpsv16_projector_closure_canonical_form_complement_expansion}\\
        &=(I-P)A^i(I-P).
    \label{eq:cpsv16_projector_closure_canonical_form_complement_invariance}
\end{align}
Applying
(\ref{eq:cpsv16_projector_closure_canonical_form_source_coinvariant_projection})
to $I-P$ gives
\begin{align}
    A^i(I-P)
        &=(I-P)A^i(I-P).
    \label{eq:cpsv16_projector_closure_canonical_form_complement_coinvariance}
\end{align}
Both off-diagonal corners therefore vanish, so
\begin{align}
    PA^i
        &=A^iP,
        \qquad\text{for every }i.
    \label{eq:cpsv16_projector_closure_canonical_form_reducing_projection}
\end{align}
Every projection used in the invariant-subspace decomposition is reducing.
This implication is formalized without a left-canonical or unital
normalization.\footnote{Formal declarations:
\leanid{MPSTensor.commutes_of_hasInvariantProjectorClosure_of_lowerZero} and
\leanid{MPSTensor.exists_reducing_projection_of_hasInvariantProj} in
\doclink{TNLean/MPS/CanonicalForm/ProjectorClosure}.}

\section{The spectral implication}

The direct-sum decomposition produces irreducible corner tensors $B_k$ and
isometries $V_k$ satisfying
\begin{align}
    A^iV_k
        &=V_kB_k^i.
    \label{eq:cpsv16_projector_closure_canonical_form_corner_intertwining}
\end{align}
This is formalized by
\leanid{MPSTensor.exists_irreducible_blockDecomp_with_isometry_of_hasInvariantProjectorClosure}
in
\doclink{TNLean/MPS/CanonicalForm/ProjectorClosureDecomposition}.

Two normalization-free spectral steps complete the source proposition. First,
each nonzero irreducible corner has a Perron eigenpair
\begin{align}
    \mathcal E_{B_k}(\rho_k)
        &=r_k\rho_k,
    \label{eq:cpsv16_projector_closure_canonical_form_perron_eigenpair}\\
    \rho_k
        &>0,
    \label{eq:cpsv16_projector_closure_canonical_form_positive_perron_vector}\\
    r_k
        &>0.
    \label{eq:cpsv16_projector_closure_canonical_form_positive_spectral_radius}
\end{align}
Rescaling $B_k$ by $r_k^{-1/2}$ moves its transfer map to spectral radius one,
while the removed scale becomes
\begin{align}
    \mu_k
        &=r_k^{1/2}.
    \label{eq:cpsv16_projector_closure_canonical_form_block_weight}
\end{align}
Second, the absence of nontrivial $p$-periodic vectors on every restriction of
$A$ to a minimal invariant subspace forces
\begin{align}
    \operatorname{spec}_{\mathrm{per}}
        (\mathcal E_{r_k^{-1/2}B_k})
        &=\{1\}.
    \label{eq:cpsv16_projector_closure_canonical_form_trivial_peripheral_spectrum}
\end{align}
Each rescaled corner is therefore normal. Both steps are formalized without a
normalization hypothesis.\footnote{Formal declarations:
\leanid{MPSTensor.isNormalTensor_invSqrt_smul_of_unique_peripheral} and
\leanid{MPSTensor.exists_normalTensor_blockDecomp_sameMPV₂Pos_of_hasInvariantProjectorClosure}
in \doclink{TNLean/MPS/CanonicalForm/ProjectorClosureSpectral}.}

The cited proposition does not assume a left-canonical identity such as
\begin{align}
    \sum_i(A^i)^\dagger A^i
        &=\Id.
    \label{eq:cpsv16_projector_closure_canonical_form_unassumed_normalization}
\end{align}
The formal completion does not assume this identity either.

\section{Zero blocks and the length-zero coefficient}

A corner whose matrices all vanish cannot be rescaled to spectral radius one.
These are the zero blocks allowed by the source convention
\begin{align}
    \sum_kD_k
        &\leq D,
    \label{eq:cpsv16_projector_closure_canonical_form_dimension_bound}
\end{align}
stated immediately after the invariant-subspace decomposition
\cite[line~219]{Cirac2016MPDO_arXiv}. Removing a zero block changes exactly one
datum of the generated family: the length-zero coefficient, which equals the
bond dimension.

The formal sufficient condition therefore concludes equality of the generated
matrix product vectors at every positive length, together with
(\ref{eq:cpsv16_projector_closure_canonical_form_dimension_bound}). A former
all-length variant assumed that no restriction of $A$ to a minimal invariant
subspace was the zero tensor. That assumption is absent from the source and
served only to recover the empty-word coefficient, so the variant has been
removed. The positive-length statement is the source-faithful treatment of
the length-zero bookkeeping analyzed in
\path{docs/paper-gaps/cpsv16_zero_tail_length_zero_decomposition.tex}.

\section{Conclusion}

The sufficient condition stated at lines~253--254 of
Ref.~\cite{Cirac2016MPDO_arXiv} is formalized. Projector closure and the
absence of nontrivial $p$-periodic vectors yield a decomposition of the
generated matrix product vectors into normal-tensor blocks with nonzero
weights, at every positive length, with the dimension bound
(\ref{eq:cpsv16_projector_closure_canonical_form_dimension_bound}). No separate
physical statement is retained for the length-zero coefficient.

\bibliographystyle{alpha}
\bibliography{references}

\end{document}
